% ============================================================================
% 问题二 —— 三段式：分析与准备 -> 建模与求解 -> 结果检验。
% 写作合同见 05-问题一.tex 顶部注释（四问同构，此处不重复）。
% ============================================================================
\section{问题二模型的建立与求解}

\subsection{分析与准备}
% 数据体检结论、口径落地、预处理、记号约定。

\subsection{模型的建立}
% 变量、目标、约束、方程。推导只留关键结论，过程移附录。

\subsection{模型的求解}
% 算法与实现要点（复杂度、分辨率预算、样本量及其来源）。
% 样本量必须写清怎么定的：先量单次耗时 -> 倒推可行规模 -> 对齐精度需求。

\subsection{结果与检验}
% 主结果（数字走宏）+ 关键图 + 证书（置信下界 / 收敛 / 双路互证）。
